\section{结果讨论}
表\ref{tab:summary}汇总了五个问题的核心结果。可以看出，这组题目并不是五个彼此独立的数值任务，而是同一条几何主线上的连续推进：Q1负责把“运动”变成可计算的参数递推，Q2把“是否碰撞”变成矩形交叠判定，Q3把“最小螺距”变成角点到板凳轴线距离的临界搜索，Q4把“调头路径”变成固定交点下的两圆弧几何构造，Q5则把“速度上限”变成比例缩放下的约束反推。

\begin{table}[H]
\centering
\caption{五个问题的核心结果汇总}
\label{tab:summary}
\begin{tabularx}{\textwidth}{>{\centering\arraybackslash}p{0.12\textwidth}L}
\toprule
问题 & 核心结果 \\
\midrule
Q1 & 0--300 s全把手轨迹生成完成，最大相邻把手距离误差为$8.32\times10^{-12}$ m \\
Q2 & 无碰撞终止时刻为$412.474$ s，首个碰撞板对为$[0,8]$ \\
Q3 & 最小螺距为$0.45$ m，粗扫触发螺距为$d=0.45$ m \\
Q4 & 固定交点两圆弧调头路径长度为$13.621$ m \\
Q5 & 固定交点路径下最大恒定速度为$1.406$ m/s \\
\bottomrule
\end{tabularx}
\end{table}

从结果稳定性看，Q1、Q2、Q3和Q4都具有明确的边界检验或区间逼近结构，因此数值结论较为稳定。Q5建立在问题四固定路径的速度矩阵之上，其速度上界由单位头速下的最大速度放大倍数直接给出。

从方法一致性看，全文始终使用同一套把手间距和几何递推框架，这一点避免了“Q1按中心线算、Q2按实体算”造成的逻辑断裂。尤其是Q3和Q5，都直接建立在前序问题给出的几何约束之上，因此它们的结论必须与前序问题保持同一解释口径。
